\begin{minipage}[t][][t]{0.82\linewidth}
    On lance une bille $\mathrm{(S)}$, de masse $m$, verticalement vers le haut
    à l'instant $(t_{0}=0)$ avec une vitesse initiale $\vec{v}_{0}$. On étudie
    le mouvement de chute libre de la bille dans un repère ($O,\vec{k}$) lié à
    la Terre supposé galiléen (Fig.~\ref{fig:bille}). On repère la position du
    centre d'inertie $G$ de la bille à un instant $t$ dans ce repère par
    l'ordonnée $z_{G}$.

\begin{enumerate}
    \item En appliquant la deuxième loi de \textsc{Newton}, montrer que
          l'équation différentielle vérifiée par l'ordonnée $z_{G}$ s'écrit~:
          $\frac{\mathrm{d}^{2} z_{G}}{\mathrm{d} t^{2}}=-g$.
    \item Quelle est la nature du mouvement de $G$ au cours de la phase de
          montée~? Justifier.
    \item L'équation horaire du mouvement de $G$ est~:
          $z_{G}=-5 t^{2}+2 t+1,5 \quad (\unit{\m})$
          \begin{enumerate}
              \item Déterminer les valeurs de $z_{0}$ et $v_{0}$ à $t_{0}=0$.
              \item À quel instant la vitesse de $G$ s'annule-t-elle~?
          \end{enumerate}
\end{enumerate}
\end{minipage}
\hfill
\begin{minipage}[t][][t]{0.16\linewidth}
\begin{figure}[H]
    \centering
    \begin{tikzpicture}
        \draw (-.5,0) -- (.5,0);
        \foreach \x in {-.5,-.4,...,.5} \draw[thick] (\x,0) -- +(-60:.3);
        \draw[-Latex,thick] (0,0)
            node[point,label=above left:$O$] {} -- ++(0,5) node[above] {$z$};
        \draw[->,very thick] (0,0) -- ++(0,1) node[right,pos=0.8] {$\vec{k}$};
        \node[circle,ball color=lightgray,draw,inner sep=9pt,
            label=right:$\mathrm{(S)}$] (S) at (1,2) {};
        \node[circle,inner sep=1pt,fill,
            label={[left=-0.5mm]:$G$}] (G) at (S.center) {};
        \draw[->,line width=0.6pt]
            (G) -- ++(0,1) node[right,pos=0.85] {$\vec{v}$};
    \end{tikzpicture}
    \caption{}
    \label{fig:bille}
\end{figure}
\end{minipage}

